\section{Lean formalization plan}\label{sec:formalization}

\subsection{Two distinct Lean layers}

A faithful formalization should not confuse a theorem in Lean's ambient
classical mathematics with a derivation in the object theory \(\ipra\).
We therefore propose two linked libraries.

\begin{enumerate}
\item \textbf{Mathlib/model layer.}  Define ordinary finite coefficients and
analytic sums over \(\Q,\R,\C\); prove the algebraic identities, convergence
bounds, and correctness of certificate generators.  This layer is useful for
rapid development and for checking formulas.
\item \textbf{Object-theory layer.}  Represent every primitive recursive
coefficient, partial sum, cutoff, and error test by a symbol/term in
\code{Prf PRAAx}; derive the first-order sentences using the checked iPRA core.
Only the second layer establishes the foundational claim of this report.
\end{enumerate}

The repository already contains a Mathlib-side definition
\code{Fabius.classicalHypergeometricTerm} for \({}_{s+1}F_s\), using
Mathlib's rising-factorial polynomial, and a noncomputable infinite sum
\code{classicalHypergeometric}.  It also proves nonvanishing and termination
lemmas in the context of a classical \(q\to1\) limit.  These definitions can
be reused for the model layer, but the noncomputable \(\sum'\) is not itself
an iPRA construction.

The checked \code{IPRADerivations} library supplies the object-theory
infrastructure needed here: integer differences, a ring tactic, definitions
by primitive recursion, finite sums and products, powers, factorials, order,
bounded search, pairing, and multi-sums.  This sharply reduces the amount of
new foundational code.

\subsection{Proposed module graph}

\begin{center}
\small
\begin{tabularx}{\textwidth}{@{}p{.28\textwidth}X@{}}
\toprule
Module & Main declarations \\
\midrule
\path{Hypergeometric/Codes} & signed rationals, complex-rational pairs,
normalized or cross-multiplied equality, one-norm, dyadic error predicate \\
\path{Hypergeometric/Pochhammer} & rising factorial numerator/denominator,
split and shift laws, negative-integer termination \\
\path{Hypergeometric/Coefficient} & \(A_n\), guard predicate, ratio
recurrence, generalized \({}_pF_q\) term \\
\path{Hypergeometric/FiniteSeries} & partial sums, Cauchy products,
derivatives, substitution in \(R/(z^N)\) \\
\path{Hypergeometric/ODE} & coefficient form of \(\mathcal L\), normalized
uniqueness, Frobenius gauge identities, Q-form \\
\path{Hypergeometric/Transform} & Euler, two Pfaff, quadratic/Kummer
certificates, six/24-form generator \\
\path{Hypergeometric/Special} & logarithm, arcsine, cubic equation, Glasser
family, Jacobi/Legendre/Chebyshev \\
\path{Hypergeometric/Bounds} & explicit ratio cutoff, geometric tail,
head--tail combinators, derivative moduli \\
\path{Hypergeometric/Integral} & endpoint cutoffs, Riemann sums, beta shifts,
Euler integral, substitutions \\
\path{Hypergeometric/Summation} & Vandermonde, Gauss chamber, Kummer,
Gauss second, Bailey \\
\path{Hypergeometric/Contiguous} & seven master identities, shift reducer,
finite Gauss fraction \\
\path{Hypergeometric/Continuation} & local Taylor recurrence, disk-step
majorant, finite path continuation \\
\path{Hypergeometric/Checks/Axioms} & \code{\#print axioms} for every exported
object-theory theorem; no \code{sorry} \\
\bottomrule
\end{tabularx}
\end{center}

The exact finite core should be implemented first.  It is independent of the
analytic-name infrastructure and will already formalize a substantial fraction
of both reference pages.

\subsection{Executable model definitions}

The companion file \code{lean/HypergeometricFinite.lean} uses a direct
rational definition suitable for experimentation:
\begin{lstlisting}[language={}]
def pochhammer (a : Q) : Nat -> Q
| 0     => 1
| n + 1 => pochhammer a n * (a + n)

def hgCoeff (a b c : Q) (n : Nat) : Q :=
  pochhammer a n * pochhammer b n /
    (pochhammer c n * n.factorial)

def hgPartial (a b c z : Q) (N : Nat) : Q :=
  sumRange N (fun n => hgCoeff a b c n * z^n)
\end{lstlisting}
The production version should use a guarded rational structure or state
nonzero hypotheses, rather than relying on Lean's total field division at a
zero denominator.

At the Mathlib level, specialization of the existing generic definition is:
\begin{lstlisting}[language={}]
def gaussTerm (a b c z : K) (n : Nat) : K :=
  Fabius.classicalHypergeometricTerm
    (s := 1) (fun | 0 => a | 1 => b) (fun _ => c) z n
\end{lstlisting}
One then proves agreement with \code{hgCoeff} over \(\Q\) and uses the
existing lemmas about \code{ascPochhammer}.

\subsection{The object-theory representation}

Working directly with rational quotients makes cancellation side conditions
awkward.  A robust object-theory implementation stores a rational expression
as a pair of signed-integer terms \((N,D)\), with \(D>0\).  Equality is
\[
  N_1D_2=N_2D_1
\]
in the ring of integer differences.  Addition and multiplication are
polynomial; inversion swaps numerator and denominator under a nonzero guard.
For most hypergeometric identities it is even cheaper to keep the coefficient
as the unreduced pair
\[
 \operatorname{num}(A_n)=\poch a n\poch b n,
 \qquad
 \operatorname{den}(A_n)=\poch c n\,n!,
\]
where each rational parameter has first been expanded into a common integer
denominator.  The core recurrence is then one integer polynomial equality.

A typical target theorem has the shape
\begin{lstlisting}[language={}]
theorem coeff_recurrence :
  ([] : List (Fm PTm 0)) |- [PRAAx]
    allN 4 (guardCoeff ==>.
      ratEq
        ((n+1) * (n+c) * coeff (n+1))
        ((n+a) * (n+b) * coeff n)) := by
  -- unfold primitive recursions; use product equations and ipra_ring
\end{lstlisting}
This is schematic notation; the archive's Lean outline records the intended
binder order and imports, not a claim that these exact identifiers already
exist.

\subsection{Induction invariants for the finite core}

The crucial proofs have unusually small invariants.

\paragraph{Pochhammer split.}
For fixed \(x,m\), induct on \(n\) with the literal equation
\[
 \poch x{m+n}=\poch x m\poch{x+m}n.
\]

\paragraph{Termination.}
Induct on \(n\le m\) for the factorial formula, then use the zero factor at
\(m+1\) and the successor equation for later indices.

\paragraph{Derivative order.}
Induct on \(r\) with a universally quantified observed coefficient \(n\), or
encode the bounded conjunction of all coefficients below \(N\).  The step is
\(\poch a{r+1}=\poch a r(a+r)\).

\paragraph{ODE.}
No induction is needed beyond the coefficient definitions: one
\code{ipra\_ring} invocation proves the cross-multiplied recurrence at an
arbitrary \(n\).

\paragraph{Vandermonde.}
Use the induction in \cref{thm:vandermonde}.  The induction hypothesis is
universal in the rational parameter codes, so it may be instantiated at
\((b+1,c+1)\) in the successor step.

\paragraph{Transformations.}
Implement a small symbolic differential-expression datatype.  A certificate
contains two rational-polynomial residuals and identities expressing them as
zero.  Reflection into the object theory turns each residual check into a
ring theorem.  Formal uniqueness then transports the result to coefficients.

\subsection{Formalizing the analytic contracts}

A dyadic complex-rational approximation can be stored as four natural
numbers: signed numerators for real and imaginary parts and a common power-of
two denominator.  Define
\[
 \operatorname{errLE}(X,Y,p)
\]
to mean, after scaling to a common denominator,
\(\norm{X-Y}\le2^{-p}\).  This is a quantifier-free formula over natural
arithmetic.

For a fixed evaluator \(E(p)\), the Cauchy contract is the sentence
\[
 \forall p\forall q\;
 (q\ge p+s\to\operatorname{errLE}(E(q),E(p+s),p)).
\]
The analytic equality of fixed evaluators \(E_1,E_2\) is
\[
 \forall p\;
 \operatorname{errLE}(E_1(p+s),E_2(p+s),p).
\]
The quantifiers are ordinary first-order quantifiers; only the evaluator
symbols are fixed by the theorem schema.

The tail theorem can be compiled without a general theory of real limits.
Define the primitive recursive cutoff \(N(a,b,c,r,p)\) by the explicit formula
and bounded search of \cref{lem:ratio-cutoff,thm:disk-convergence}.  Prove
inside iPRA that every finite tail segment between \(N\) and \(M\) is bounded
by the geometric expression.  The Cauchy theorem follows directly.  This
approach is shorter than first building a full abstract constructive analysis
library.

\subsection{Formalizing Euler's integral}

A minimal development needs only the following reusable lemmas:
\begin{enumerate}
\item a bounded-sum estimate for rational complex numbers;
\item a modulus-to-Riemann-sum Cauchy lemma on a closed rational interval;
\item endpoint power-tail bounds with a positive rational exponent margin;
\item summation/integration interchange for a finite head plus a uniform
remainder;
\item finite substitution for the specific rational maps used in this report.
\end{enumerate}
Each theorem is a first-order statement about explicit sums and inequalities.
The integral is not introduced as a higher-order object; every named integrand
gets its own primitive recursive Riemann-sum symbol.

\subsection{Suggested implementation milestones}

\begin{enumerate}[label=\textbf{M\arabic*.}]
\item Compile Pochhammer, coefficients, termination, derivatives, ODE,
formal uniqueness, and Vandermonde.  This should contain no analytic names.
\item Add a finite truncated-series evaluator and reflection checker.  Prove
Euler, Pfaff, the printed quadratic transformations, the cubic equation, and
the Glasser family.
\item Add rational/complex-rational bounds and prove the strict-disk tail
modulus and termwise differentiation.
\item Add beta Riemann sums and prove the Euler-chamber integral and Gauss
summation; reuse the gamma packet for the quotient form.
\item Add Kummer, Gauss second, Bailey, incomplete beta, and elliptic
specializations.
\item Add path stepping and conditional certificates for the connection
formula, continued fraction, and Barnes contour.
\end{enumerate}

\subsection{Expected theorem declarations}

The final object library should export, at minimum, declarations corresponding
to the following mathematical statements:
\begin{lstlisting}[language={}]
poch_split, poch_shift, poch_neg, poch_zero_after
hg_coeff_ratio, hg_terminates, hg_deriv_coeff
hg_ode_coeff, hg_regular_unique
hg_euler_formal, hg_pfaff1_formal, hg_pfaff2_formal
hg_quadratic_formal, hg_kummer_quadratic_formal
hg_contiguous_master, hg_shift_reduction
chu_vandermonde
hg_cubic_algebraic, glasser_family
hg_ratio_cutoff, hg_tail_cauchy, hg_termwise_derivative
euler_integral_effective, gauss_sum_effective
kummer_sum_effective, gauss_second_effective, bailey_effective
hg_path_continue, gauss_fraction_finite
\end{lstlisting}
The last two analytic groups should state all margins and moduli in their
types; hidden typeclass assumptions such as abstract completeness should not
be used to discharge them.

\subsection{Audit status of the supplied Lean files}

The archive includes:
\begin{itemize}
\item \code{HypergeometricFinite.lean}, a detailed Mathlib-style finite-core
outline with definitions and many proof scripts or proof sketches;
\item \code{HypergeometricIPRA.lean}, the object-theory declaration plan and
formula encodings;
\item \code{FORMALIZATION\_PLAN.md}, a file-by-file dependency and proof
checklist.
\end{itemize}
The environment used to prepare this package did not contain a Lean executable,
so these files were not compiled.  They are explicitly labeled as an outline;
no claim of kernel verification is made.  The existing repository files cited
above are separate and already provide checked infrastructure.
